\documentclass[11pt,letterpaper]{article}

%\renewcommand{\familydefault}{\rmdefault}


% aas: use the margin parameter to adjust margin if too long, 1" is standard for nsf
\usepackage[margin=1.0in]{geometry}
\usepackage{wrapfig}
\usepackage{mathptmx} 
\usepackage{array}
\usepackage{calc}
\usepackage{graphicx}
\usepackage{amssymb}
\usepackage{color}
\usepackage[most]{tcolorbox} % for colored/framed boxes
\usepackage{xcolor}          % color support
\usepackage{ifthen}
\usepackage{url}
\usepackage{amssymb}
\usepackage{subcaption}
%\usepackage{fancyhdr}
\usepackage{enumitem}
\usepackage{cite}
\usepackage{caption}
%\input{defs}
\usepackage[utf8]{inputenc}
 \usepackage{tgbonum}
\usepackage[osf]{mathpazo}
\usepackage{mathrsfs, amsmath}
\usepackage{amsthm}
\usepackage{tikz}
\usepackage{pbox}
\usepackage{mdframed}
\usepackage{titlesec}
\titlespacing*{\section}{0pt}{0pt}{0pt}
\titlespacing*{\subsection}{0pt}{0pt}{0pt}
\titlespacing*{\subsubsection}{0pt}{0pt}{0pt}
\usetikzlibrary{arrows,decorations.pathmorphing,backgrounds,fit,positioning,shapes.symbols,chains, shapes,snakes}
\usepackage{bm}
\usepackage{comment}
\usepackage[normalem]{ulem}
\usepackage{multirow}
\usepackage{booktabs}
\usepackage{tabularx}
\usepackage{float}
\usepackage{adjustbox}
\usepackage{natbib}
\usepackage[colorlinks=true, linkcolor=blue, citecolor=blue]{hyperref}


\usepackage{setspace}
\setstretch{1.0} 
\usepackage{anyfontsize}
\fontsize{10}{12}\selectfont 


\usepackage{todonotes} % for comments
\newcommand{\jan}[1]{\todo[inline,color=blue!20!white]{\textbf{Jan:} #1}}

%\titleformat{\section}
%  {\large\bfseries\color{black}}
%  {\thesection}{1em}{}
  
%  \titleformat{\subsection}
 % {\normal\bfseries\color{black}}
 % {\thesubsection}{1em}{}
  
\newcommand{\tosay}{true}
\newcommand{\elza}[1]{ %
{\ifthenelse{\boolean{\tosay}}{\begin{quotation}\textcolor{red}{Elza: #1}\end{quotation}}{}}}
\newcommand{\efe}[1]{ %
{\ifthenelse{\boolean{\tosay}}{\begin{quotation}\textcolor{green}{Efe: #1}\end{quotation}}{}}}
\newcommand{\sid}[1]{ %
{\ifthenelse{\boolean{\tosay}}{\begin{quotation}\textcolor{magenta}{Siddharth: #1}\end{quotation}}{}}}
\newcommand{\mc}{\mathcal}
\newcommand{\pr}{\mathbb P}
\newcommand{\agn}[1]{
	\begin{align}
	#1
	\end{align}
}

\theoremstyle{plain}
\newtheorem{theorem}{Theorem}
\newtheorem{corollary}{Corollary}[theorem]
\newtheorem{prop}{Proposition-}

\theoremstyle{definition}
\newtheorem{defn}{Definition}


\theoremstyle{remark}	
\newtheorem{remark}{Remark}


\DeclareMathOperator*{\argmax}{arg\,max}
\DeclareMathOperator*{\argmin}{arg\,min}
\newcommand{\ex}{\mathbb E}
\newcommand{\ind}{\operatorname{\mathds{1}}}
\allowdisplaybreaks



\begin{document}


% \pagestyle{empty}
\setcounter{page}{1}



%\bibliography{nyurefs,burefs}
%\newpage\pagestyle{fancy}\fancyhf{}
%\fancyfoot[C]{\footnotesize R-\thepage}
%\setcounter{page}{1}




%\newpage
%\fancyfoot[C]{\footnotesize \thepage}
\setcounter{page}{1}
\begin{center}
{\Large \textbf{
NASA Space To Soil Challenge
}} 
\end{center}

% Judges Criteria:
% The solution proposes a novel adaptive sensing or onboard processing approach to enhance regenerative agriculture, forestry, or a similar land resilience objective.
% Solution design articulates system integration thinking and provides a full-stack implementation plan with external technologies.
% End user specifications are well-defined for application to regenerative agriculture, forestry, or a similar land resilience objective.
% The solution provides evidence that it can function as described within SmallSat constraints.
% The solution results in a measurable and significant improvement to the science need, through the use of adaptive sensing or onboard processing aligned with mission requirements.
% The potential impact of the solution to regenerative agriculture or forestry is quantifiably and convincingly articulated (such as data downlink reduction, cost efficiencies, or SWAP metrics).
% The solution applies market analysis to identify industry gaps and describe the development process from base concept to Minimum Viable Product (MVP).
% There is thorough discussion of the solution's market generalization to address applications beyond agricultural or forestry purposes.
% There is clear evidence of the team's quallifications across domains for project execution.
% The submission paper clearly articulates the solution, its benefits, and its feasibility.
% The video presentation effectively communicates the solution's benefits and key design aspects.
% The solution's key capabilities are substantiated by the provided software code (software) or schematics/diagrams (hardware).

\section{Introduction} \label{sec:intro}
% Brief overview of the proposed solution and description of your team/organization.\\
% Provide a condensed summary of how your solution addresses all 5 of the Judge Criteria concepts.

Soil is one of the most essential foundations of modern life because it quietly supports nearly every system humans depend on. It grows the food we eat, filters the water we drink, and stores vast amounts of carbon that help regulate the global climate. Healthy soil sustains biodiversity—from microbes to insects to plant roots—that keeps ecosystems functioning and resilient. It also underpins infrastructure by stabilizing landscapes, reducing erosion, and buffering floods. In a world facing climate change, population growth, and increasing pressure on natural resources, protecting soil health is not just an environmental concern but a fundamental requirement for food security, clean water, and long‑term human well‑being. Many modern day scientific practices such as regenerative/ precision agriculture practices like remote‑sensing–based yield prediction, variable‑rate fertilization, soil moisture mapping, drone‑based crop monitoring all depend on predictable soil structure and fertility. Similarly, Soil degradation reduces soil organic carbon which disrupts carbon sequestration estimates and weakens climate mitigation models, complicates MRV systems (Monitoring, Reporting, Verification) for carbon markets. Healthy soil is essential for reliable climate projections and carbon accounting. Sustainable Forestry \& Agroforestry is another area which is highly dependent on soil health. Degraded soils reduce tree growth modeling accuracy, carbon stock estimates and resilience assessments. Degraded soils produce noisy or misleading signals, reducing the accuracy of these technologies. 

Apart from the benefits that real time soil erosion monitoring can accrue to the fields of regenerative agriculture and sustainable forestry, there are a host of other land resilience applications where real time erosion monitoring can be a game changer. Some such application areas are hydropower generation protection, siltation mitigation, slope failure prediction, preventing desertification and associated humanitarian crises etc. Hence, the detection of soil erosion in real time can act as a key enabler in the areas of regenerative agriculture, regenerative forestry and a host of land resilience applications. 

This motivated us to envisage a novel application SoilSentinel, an advanced CubeSat concept designed by Darwin Geospatial and Johns Hopkins University for the 2026 NASA Space-to-Soil Challenge. This satellite aims to detect soil erosion and land degradation in real time by utilizing a distilled AI model and onboard processing to filter data. By maintaining spatial embeddings and identifying deviations from seasonal baselines, the system can autonomously trigger alerts while reducing bandwidth requirements by 95\%. Such real-time soil degradation inputs can prove to be valuable for entities like Institutional Land Investors and Real Estate Investment Trusts to protect their book value as land is an asset for them; Global Food and Beverage Corporations to protect their supply chains; Agricultural Insurance Providers to alert their insurers of possible crop failure and provide early assistance to mitigate a crisis thereby reducing payout volatility; Governments will also benefit from such information as it gives the opportunity to detect desertification at early stages which can lead to loss of arable land and thereby impact food security; Forestry companies managing timber lands for 30 to 50 year cycles can use such information to ensure post-harvest land recovery.



\section{Proposed Solution Details}
% Explanation of proposed adaptive sensing or onboard processing approach to regenerative agriculture or forestry, or a similar land resilience objective. \\
% Detail the problem statement for your solution.\\
% Discuss the selected land use algorithm(s) and small satellite technical requirements. \\
% Outline mission requirements in alignment with small satellite onboard constraints such as power, compute, and bandwidth.\\
% Explain solution design integration into a current commercial market full-stack implementation system.
The proposed solution consists of the SoilSentinel advanced 6U CubeSat concept specifically designed to detect soil erosion and land degradation autonomously and in real time. It is conceptualised to be built on a Terran Orbital TRIUMPH satellite bus weighing approximately 12 kg,  utilizing a 100 W deployable solar array and features an X-band downlink capable of 50 Mbps. To capture detailed land surface data, the satellite is equipped with dual optical sensors: a Simera MultiScape100 CIS for Visible and Near-Infrared (VNIR) imagery across 8 bands, and a Dragonfly Chameleon SWIR sensor offering 4 Shortwave Infrared bands without the need for cryogenic cooling. The core intelligence of the system is powered by an onboard NVIDIA Jetson Orin Nano processor, which executes complex AI inference at the edge within strict power and size constraints (15W and 8 GB RAM).

Instead of relying on ground-based computing, SoilSentinel processes all captured imagery directly in space using a highly optimized, 3-megabyte MobileUNet "student" model. Prior to launch, this model is created via knowledge distillation from a massive 300-million-parameter Prithvi-EO-2.0 "teacher" model trained on NASA Harmonized Landsat Sentinel (HLS) data. During a satellite overpass, the sensors capture raw 12-band imagery, which is radiometrically calibrated down to five essential bands (Blue, Green, Red, Near-Infrared, and Shortwave-Infrared 1). The system then divides the images into 224x224 pixel chips, aggressively masking out clouds and oceans to reject unusable data, before passing the remaining imagery through the MobileUNet model to generate a degradation risk score for every pixel in roughly 0.05 milliseconds per chip.

The most groundbreaking aspect of the SoilSentinel mission is its novel approach to monitoring land changes using 'persistent spatial embeddings'. Unlike traditional satellites that simply capture and transmit massive raw image files, SoilSentinel maintains a continuous mathematical memory of the Earth's surface. Specifically, it assigns and stores a '128-dimensional memory vector' for every single square kilometer of monitored cropland. This embedding acts as a localized, highly compressed digital signature of the land's current state, allowing the satellite to remember what specific parcels of land typically look like without needing constant interaction with ground stations. Presently, no existing system utilizes this kind of onboard spatial memory bank.
This persistent memory is actively managed using an Exponential Moving Average (EMA). Every time the satellite passes over a monitored region, the newly computed 128-dimensional embedding is mathematically blended with the historical baseline. These embeddings are used to derive indices such as the Normalized Difference Vegetation Index (NDVI) and the Bare Soil Index (BSI). Revised Universal Soil Loss Equation (RUSLE) is the industry standard for erosion modeling. In this method the C-factor (Cover-Management) is derived from NDVI. Research has shown that there exists an exponential relationship between NDVI and the C-factor, demonstrating that as NDVI increases, the risk of water-induced soil erosion drops exponentially \cite{van2000soil}. NDVI below 0.3 to 0.4 typically indicates a "tipping point" where the soil surface is no longer protected by a canopy, leading to a massive increase in splash erosion. While NDVI tells us where the vegetation is, the Bare Soil Index (BSI) is used to characterize the soil's physical state. Another work, \cite{diek2017barest} utilized BSI combined with SWIR bands to differentiate between healthy topsoil and exposed, degraded subsoil. Remote sensing inputs can be used to formulate a Soil Degradation Index (SDI) \cite{nascimento2021soil} that can directly be used to assess the present state of the soil. The SDI will enable to map and assess degradation severity across large areas in several ways. The SDI can be used to group areas into categories of degradation risk like very low, low, medium, high, and very high.
% The placement of an area into one of these severity levels is primarily guided by crucial soil attributes, namely clay content and Cation Exchange Capacity (CEC). Locations that possess greater amounts of these structural attributes are classified at a lower risk of degradation.

The SDI quantifies severity by linking satellite imagery directly to physical soil health. As the SDI severity level increases from "low" to "high," the spectral reflectance of the bare soil also increases. This elevated reflectance is a direct indicator of a loss of clay and soil organic matter (OM), which characterizes highly degraded, sandy areas. 
% To confirm the accuracy of these severity levels, the SDI is validated against the amount of OM present. Areas classified by the SDI as having the highest severity of degradation consistently generally correspond with the lowest OM content.
By assigning these severity values to specific geographic coordinates, the SDI is capable of creates continuous, easily understood maps of soil conditions. For example, when researchers applied the SDI to a study region in São Paulo, Brazil, it was able to precisely quantify that 15\% of the total area suffered from an extreme (very high) degradation risk, while 24\% of the land was at a very low risk. Another work studied on how remote sensing indices such as Normalized Difference Vegetation Index (NDVI), Normalized Difference Water Index (NDWI), Soil Moisture Index (SMI), and Leaf Area Index (LAI) can be linked to soil degradation \cite{chandramohan2024evaluating}. The indices can be used to mathematically quantify the exact severity of the degradation using derived Land Degradation Vulnerability Index (LDVI). The LDVI can be used to classify land into four categories based on their gradation of vulnerability to land degradation: (a) nonaffected zones (N) with LDVI value $\leq$ 1.17; (b) potential vegetative cover (P) with LDVI value $1.17 - 1.22$; (c) fragile zones with LDVI value $1.22 - 1.37$; and (d) critical zones LDVI value is $\geq 1.37$.

% Regression Models with Surface Albedo Because land degradation causes a decrease in vegetation and an increase in bare, bright soil, researchers can quantify degradation severity by calculating the linear regression between surface albedo and vegetation indices like NDVI. By extracting the mathematical slope coefficient between Albedo and NDVI across thousands of pixels, researchers generate an equation to score the degradation intensity of any given area. This allows a landscape to be categorized into exact severity classes (e.g., from "null" to "very high" degradation risk). For example, when applying an Albedo-NDVI regression model to a study site in Iran, researchers could precisely quantify that 44.32\% of the land was in a state of "very high" degradation. GIS-Based Vulnerability Models (e.g., MEDALUS) use indices like NDVI, NDWI, SMI, and LAI often combined in a Geographic Information System (GIS) to calculate a Land Degradation Vulnerability Index (LDVI).
% Researchers use spatial analysis tools to extract pixel values for all four of these indices simultaneously. By assessing the combined values, the model categorizes the land into non-affected zones, potential vegetative cover, fragile zones, and critical zones.
% This quantification process consistently proves that areas falling into the "critical" degradation zones are mathematically defined by significantly lower NDVI, LAI, and NDWI values, alongside distinct SMI ranges that point to severe moisture deficits and biomass loss. Conversely, non-affected areas maintain high index values indicative of healthy, moist, and well-covered soil.

Another study \cite{ebrahimi2024assessing} used a simplistic BSI-NDVI model to categorize land degradation. The model is built upon the strong negative correlation between BSI and NDVI. As bare soil increases in a degraded area, vegetation cover naturally decreases. The study establishes a linear regression between these two indices to generate a specific mathematical equation for degradation intensity. Specifically, the land degradation severity (I) for this model is calculated using the equation: $I = 1.343 \times BSI - NDVI$. The numerical values generated by the BSI-NDVI equation are grouped using statistical classification methods (such as the Jenks natural breaks classification) to divide the landscape into five distinct levels of degradation severity: Null, Low, Medium, High, and Very High. By applying this equation across satellite imagery (like Sentinel-2 data), researchers can map the exact percentage of land falling into each severity class. For example, when applying the NDVI-BSI model to a study site in the Rutak region of Iran, researchers quantified that 37.81\% of the total area was in a state of "high" degradation, while only 0.58\% was classified as having "null" (no) degradation. The BSI-NDVI model is considered a highly accurate tool for quantifying degradation. When researchers validated the model against actual field observations using statistical tests (such as the Mann-Whitney U-Test), they found that the severity levels predicted by the BSI-NDVI model did not significantly differ from the ground-truth field data. The model achieved a high regression coefficient of 0.86, demonstrating its reliability for estimating land degradation severity without requiring extensive physical field surveys.

The onboard AI continuously compares the real-time embedding against monthly baseline expectations uploaded from NASA MODIS datasets. The system specifically watches for significant "drifts" from these baselines, such as a drop in the NDVI greater than 0.15, or a rise in the Bare Soil Index (BSI) greater than 0.10 or even a composite metrics like the SDI, LDVI or a BSI-NDVI composite value as mentioned above. These metrics serve as direct indicators of anomalous land surface changes, such as sudden crop failure, ongoing desertification, or topsoil loss. When an embedding drifts beyond its normal seasonal parameters, the satellite acts autonomously. Rather than downlinking a massive 200 MB raw scene to Earth, the system triggers a targeted alert containing coordinates, a severity score, and an adaptive re-imaging schedule. It downlinks only these tiny ~50 KB anomaly alerts, effectively reducing the satellite's bandwidth usage by approximately 95\%. Normal, unflagged data is simply discarded onboard. Furthermore, the system autonomously instructs the satellite to adaptively re-image the flagged problem areas on its very next orbit, creating a dynamic, highly responsive sentinel for global land degradation.

\section{Potential Impact}
% Benefits and impact of the proposed solution.\\
% How does this solution build upon current research and market capabilities?\\
% What effect does the solution have on technical variables (as defined in Proposed Solution Detail) compared to traditional methods for regenerative agriculture or forestry?\\
% Provide concrete examples and scenarios where the solution outperforms current adaptive sensing or onboard processing capabilities, and highlight solution limitations.\\

Land degradation is an underlying problem that acts as precursor to many natural calamities and leads to cascading environmental effects. It has both short term effects such as landslides and long term effects such as desertification. As such, detecting the early signs of land degradation can be an enabling factor to prevent or reduce many environmental problems. It is multi-faceted and can manifest itself through a plethora of environmental problems. As such, early detection of land degradation will have a potential impact in mitigating or reducing the severity of many environmental problems. Some crucial applications where real time land degradation can have significant impacts are enumerated in more detail in the succeeding paragraphs.

Regenerative agriculture is the process of restoring degraded soils using management practices (e.g., adaptive grazing, no-till planting, no or limited use of pesticides and synthetic fertilizer, etc.) based on ecological principles. Regenerative agriculture strives to work with nature rather than against it. It is more than just being sustainable. It is about reversing degradation and building up the soil to make it healthier than its current state. Therefore, real time  detection of land degradation can act as the "nervous system" of a farm—providing the immediate feedback needed to steer complex ecological processes and determine if a certain process is proving to be effective. Regenerative practices, such as intensive rotational grazing or multi-species cover cropping, are dynamic. Traditional soil testing is far too slow to inform daily operations. For example, Real-time monitoring can detect when "trampling" or "overgrazing" starts to expose bare soil. This allows a farmer to move livestock immediately before the protective soil crust is broken and erosion begins. Another example where such real time detection would be extremely helpful is satellite data detecting a "thinning" of the biomass or a structural failure in the soil after a minor rain event. Managers can intervene with targeted seeding or mulching before the next major storm. For large-scale food companies (the Corporate Ag consumer base), real-time detection in their supply chain provides an early warning system for regional crop failures, allowing them to shift procurement strategies before a local land collapse impacts their bottom line. Regenerative agriculture is increasingly funded by Carbon Credits and Ecosystem Service Payments. These markets demand high-fidelity verification (MRV: Monitoring, Reporting, and Verification). Soil organic carbon is highly correlated with soil structure and stability. Real-time erosion detection proves that the "sequestered" carbon isn't being washed away into local streams, providing the "permanence" guarantee that carbon buyers require. In precision agriculture, which is the focus area for agri-tech, instead of treating the whole field as a single unit, technology allows for surgical precision. Some examples are nutrient application can be adjusted based on real-time runoff risk to prevent leaching, tilling/ disturbance can be restricted in areas where real-time sensors show high structural fragility, dynamic irrigation can be employed that pauses when sensors detect initial signs of soil slumping or soil ponding.

In sustainable forestry, the management cycle spans decades, but the most catastrophic soil loss often occurs in a matter of hours, typically during timber harvesting, road construction, or immediately following a wildfire. Real-time detection can track the regrowth of the herbaceous layer. If a specific patch isn't greening up or shows structural settling, managers can deploy targeted "hydroseeding" or mulch to stabilize the area before the first winter storms. Wildfires create "hydrophobic" (water-repellent) soils. When rain hits these areas, it doesn't soak in; it slides, leading to debris flows. In burnt-over areas, real-time detection acts as a life-safety system for downstream communities. By identifying the exact moment of soil failure on a charred ridge, the system provides a lead time for evacuations or emergency drainage work. Instead of expensive blanket aerial seeding after a fire, real-time data identifies the 5\% of the landscape that is at the highest risk for mass erosion, allowing for a 90\% reduction in mitigation costs. Real-time monitoring provides an immutable digital audit for carbon and biodiversity credits. It proves to investors that the timber was extracted with Zero Net Erosion.

Protecting hydropower and water management is another area where real time erosion detection can help prevent siltation crisis where the washed top soil ends up in reservoirs downstream causing siltation and posing a threat to expensive hydroelectric turbines via abrasion. By detecting a drift in land stability upstream, water authorities can deploy temporary sediment traps or adjust dam operations before a sediment plume hits the intake. This extends the operational life of multi-billion dollar energy assets. Land degradation is often the precursor to catastrophic mass-wasting events like landslides. SoilSentinel’s persistent embeddings can detect the subtle slumping or loss of vegetation armor on steep slopes that overlook highways, rail lines, or pipelines. Instead of reacting to a collapsed tunnel or a ruptured gas line, civil engineers can be given an autonomous alert of structural drift, allowing for pre-emptive stabilization. Real-time detection to identify degradation hotspots in remote regions can enable authorities to deploy measures like assisted natural regeneration to stabilize a community's food source before environmental failure triggers a migration crisis.


\section{Team Experience and Market Evaluation}
% Qualifications and expertise of team members (multidisciplinary teams are encouraged).\\
% Past projects or research relevant to the proposed solution.\\
% Analysis for how the proposed solution can eventually be scaled or integrated into future markets.\\
% Identification of potential product users and how target users inform solution design criteria.\\
% Analysis of potential market applications beyond agricultural or forestry purposes.\\
Executing the SoilSentinel project requires a highly specialized, multidisciplinary team of experts with skillsets spanning artificial intelligence, earth sciences, mathematical validation, and aerospace systems engineering. Primarily, the project demands AI and Machine Learning Engineers to design the complex data pipelines and handle model training. These experts must possess deep knowledge of neural network optimization, specifically utilizing knowledge distillation and INT8 quantization to shrink a massive 300-million-parameter teacher model down to a highly efficient 3-megabyte student model. Furthermore, the mission relies heavily on Geospatial, Earth Science and Remote Sensing Experts capable of processing complex satellite imagery and validate the algorithm results. Because the satellite autonomously monitors mathematical drifts in persistent spatial embeddings, the team also requires Mathematical and Algorithm Validators to rigorously review the distillation loss functions, exponential moving average (EMA) logic, and anomaly detection equations to ensure absolute numerical consistency. In addition, Cloud and Software Engineers are essential for building the underlying computing infrastructure, which includes managing GPU spot instances, CI/CD pipelines, and experiment tracking on platforms like Google Cloud to handle the intense pre-launch training loads. 
% Finally, successfully deploying this technology in space requires Space Systems and Hardware Expertise to meticulously manage Size, Weight, and Power (SWaP) budgets. These systems experts validate power consumption, optimize the 50 Mbps X-band downlink capabilities, and ensure the onboard NVIDIA Jetson Orin Nano processor safely executes inference within the strict 15W power constraints of the 12 kg, 6U CubeSat framework.

Our multi-disciplinary team consists of members that possess skillsets across the domains mentioned above. A short description of the team members along with some of their relevant work is mentioned below:
\begin{itemize}
\item Ben is a Professor and the chair of the Department of Earth \& Planetary Sciences at Johns Hopkins University. He is the scietific lead for this project. He is a globally recognized authority in hydroclimatology and has extensive experience in translating satellite-derived data into actionable food security and water management solutions. As a leader in Earth and planetary sciences, he brings a deep understanding of the physical processes governing land-atmosphere interactions, which is essential for grounding predictive land models in climatological reality. His long-standing involvement with large-scale monitoring initiatives, such as those related to famine early warning and regional climate resilience, provides the strategic vision needed to ensure that remote sensing technologies deliver tangible socio-economic benefits. Ben’s ability to bridge the gap between complex climate modeling and practical applications for humanitarian and environmental stability makes him an ideal lead for scientific validation and high-level strategic impact.

\item Gabriel is a the lead project member and is resposible for overall steering of the project. He has a rare combination of deep domain expertise in the geosciences and a sophisticated command of machine learning for environmental monitoring. His research background at the Scripps Institution of Oceanography demonstrates a proven ability to process complex, high-dimensional geospatial datasets using specialized tools like xarray and PyTorch. Gabriel excels at bridging the gap between theoretical data science and the practical physical constraints of Earth observation, making him uniquely qualified to develop and refine the AI pipelines required for real-time land and ocean resilience. His experience in translating raw sensor data into actionable ecological insights ensures that the project remains scientifically rigorous while maintaining the high-performance standards necessary for edge-computing and satellite-based applications. he also an expert in cloud computing and has deep knowledge of the Google Cloud Platform. He will thus be the key member for ensuring project scalability and deployment on the cloud.

\item Victor is the machine learning lead member. He has exceptional expertise at the intersection of machine learning and signal processing. He has demonstrated a unique ability to bridge complex mathematical theory with practical, robust architectural designs, particularly in the realm of Graph Neural Networks (GNNs). His innovative research on stabilizing deep learning filters through neighborhood-based methods shows a commitment to creating reliable, scalable AI solutions that overcome fundamental numerical challenges in high-dimensional data. Beyond his technical proficiency, Victor’s contributions to peer-reviewed publications and open-source communities reflect a collaborative spirit and a drive for academic and engineering excellence. He brings a rare combination of theoretical depth and high-performance problem-solving skills that make him an asset to any cutting-edge research or development team.

\item Sandupal is a PhD candidate at Johns Hopkins University. His research lies at the intersection of machine learning and atmosphere/ ocean science. In addition, he has extensive experience as a operational meteorologist/ oceanographer with the Indian Navy. His expertise in applying machine learning to earth science problems will be  valuable across multiple aspects of the projects. He has the necessary knowledge to bridge the gap between machine learning and earth sciences, which is helpful in interpreting the predictions given by machine learning model and identifying areas of improvement. He is also the lead scholar as he has substantial experience in scientific writing and report writing.
\end{itemize}



\section{Expected Outcomes} 
% Articulate your development process from base concept to Minimum Viable Product (MVP) based on market analysis and industry gaps.
% What are your next steps if selected as a Finalist?
The aim would be to put in place a Minimum Viable Product (MVP) for the agriculture industry using the SoilSentinel concept as the agro industry is one of the largest entities most severely affected by land degradation as it is their most important aspect. Developing the MVP would consist of sequence of steps as envisaged below:
\begin{itemize}
    \item Stakeholder requirements: Interview agro firms to determine their exact requirements and ideal response time fo the system that would render maximum benefits.
    \item Model Optimization: Convert the 300M parameter Prithvi model into a lightweight 3MB MobileUNet and use TensorRT for hardware acceleration on the Orin Nano.
    \item Metric Calibration: Establish the "Drift Thresholds." Define what constitutes a "Critical Alert" for agro-stakeholders (e.g., a specific BSI-NDVI/ SDI/LDVI delta that correlates to 1cm of topsoil loss). This phase can consist of logic encoding like programming the AI to compute Land degradation Severity (I) $I = 1.343 \times BSI - NDVI$. "Tipping Point" trigger can also be incorporated so that the system will flag any pixel where NDVI drops below 0.3–0.4, triggering an immediate "Splash Erosion" risk alert. This stage may also consist of incorporating workflows like training the "Student" model to recognize the four LDVI (Land Degradation Vulnerability Index) zones using the thresholds: Nonaffected: $LDVI \leq 1.17$, Potential: $1.17 - 1.22$, Fragile: $1.22 - 1.37$, Critical: $\geq 1.37$.
    \item Synthetic Environment Testing: Run the inference pipeline using historical NASA HLS data to simulate satellite overpasses. Validate the 0.05ms per chip processing speed.
    \item Site Selection: Partner with an institutional land investor or a regenerative ag firm to monitor a high-risk "Alpha Site" (e.g., a 10,000-hectare cattle ranch or soybean farm).
    \item Digital Proxy: Use existing Sentinel-2 or Planet imagery to "mimic" the SoilSentinel pipeline. Generate the 50 KB alerts and deliver them to the farm manager's dashboard.
    \item Feedback Loop: Interview the stakeholders. Does the alert arrive in time for them to change their tilling schedule or move livestock? Refine the "Decision Metrics" based on these operational realities.
    \item FlatSat Testing: Integrate the Jetson Orin Nano with the Simera VNIR and Dragonfly SWIR sensors on a "FlatSat" (a tabletop version of the satellite).
    \item Onboard Memory Loading: Upload the "Persistent Spatial Embeddings" for the Alpha Site. Ensure the 128D vectors for every square kilometer are correctly stored and indexed.
    \item Communication Bridge: Test the X-band downlink logic. Verify that the system can autonomously discard "Normal" data and prioritize the 50 KB "Anomaly Alerts."
    \item Orbital Commissioning: Launch the 6U CubeSat and perform a 30-day health check of the solar arrays and propulsion.
    \item Autonomous Alert Generation: During the first overpass of the Alpha Site, the satellite must capture 12-band imagery, perform onboard inference, identify a "Drift" vs. the monthly baseline and downlink the 50 KB alert via X-band.
    \item Ground-Truth Audit: Cross-reference the satellite alert with on-the-ground soil sensors or drone flights to verify accuracy.
\end{itemize}

 
\section{Conclusion}

\bibliographystyle{ieeetr}
\bibliography{soil}


\end{document}


